\section{Entscheidungsbäume}
\setcounter{aufgabe}{0}


% ============================================================
% AUFGABE 1 (einfach)
% ============================================================
\aufgabe{
Gegeben sei der folgende Trainingsdatensatz. Zielattribut ist \textbf{Grillen}.
\begin{center}
{\scriptsize
\renewcommand{\arraystretch}{1.25}
\setlength{\tabcolsep}{6pt}
\rowcolors{2}{SecondaryColor!12}{white}
\begin{tabular}{c c c c}
\rowcolor{PrimaryColor}
\textcolor{white}{\textbf{Nr.}} &
\textcolor{white}{\textbf{Wochenende}} &
\textcolor{white}{\textbf{Regen}} &
\textcolor{white}{\textbf{Grillen}}\\
1 & Ja   & Nein & Ja\\
2 & Ja   & Nein & Ja\\
3 & Ja   & Ja   & Nein\\
4 & Nein & Nein & Nein\\
5 & Nein & Nein & Nein\\
6 & Nein & Ja   & Nein\\
\end{tabular}}
\end{center}

Erstellen Sie mit dem ID3-Algorithmus einen Entscheidungsbaum:
\begin{enumerate}[label=\alph*)]
  \item Berechnen Sie die bedingte Entropie $H(V \mid W)$ für beide Attribute
        und bestimmen Sie das Wurzelattribut.
  \item Setzen Sie das Verfahren rekursiv fort und zeichnen Sie den vollständigen
        Entscheidungsbaum.
\end{enumerate}
}

\loesung{
Es genügt, die bedingte Entropie $H(V \mid W)$ zu berechnen und das
\textbf{kleinste} zu wählen, da $H(V)$ am Knoten konstant ist.

\begin{enumerate}[label=\alph*)]
  \item Von $6$ Beispielen sind $2\times$ Grillen ja und $4\times$ nein.
  \[
    H(\text{Grillen} \mid \text{Wochenende}) = -\Bigl[\,
      \underbrace{\tfrac{3}{6}}_{\text{Wo. ja}}
      \bigl(\tfrac{2}{3}\log_2\tfrac{2}{3} + \tfrac{1}{3}\log_2\tfrac{1}{3}\bigr)
      +
      \underbrace{\tfrac{3}{6}}_{\text{Wo. nein}}
      \bigl(\tfrac{3}{3}\log_2\tfrac{3}{3}\bigr)
    \,\Bigr] \approx 0{,}4592
  \]
  \[
    H(\text{Grillen} \mid \text{Regen}) = -\Bigl[\,
      \underbrace{\tfrac{4}{6}}_{\text{R. nein}}
      \bigl(\tfrac{2}{4}\log_2\tfrac{2}{4} + \tfrac{2}{4}\log_2\tfrac{2}{4}\bigr)
      +
      \underbrace{\tfrac{2}{6}}_{\text{R. ja}}
      \bigl(\tfrac{2}{2}\log_2\tfrac{2}{2}\bigr)
    \,\Bigr] \approx 0{,}6667
  \]
  Kleinste bedingte Entropie: \textbf{Wochenende} $\Rightarrow$ Wurzelattribut.

  \item Der Zweig \textbf{Wochenende = nein} ist rein (nur Grillen nein), also ein
  Blatt. Im Zweig \textbf{Wochenende = ja} (Beispiele 1--3) verbleibt nur das
  Attribut \textbf{Regen}, das die Restmenge exakt trennt:
  \[
    \text{Regen nein} \Rightarrow \text{Grillen},
    \qquad
    \text{Regen ja} \Rightarrow \text{nicht Grillen}.
  \]

  \vspace{0.4em}
  \centering
  \resizebox{0.6\linewidth}{!}{%
  \begin{tikzpicture}[
    test/.style={draw=PrimaryColor, very thick, rounded corners=4pt,
                 align=center, inner sep=5pt, text width=2.0cm, font=\footnotesize},
    rad/.style={draw=ExampleTitle, very thick, ellipse, align=center,
                inner sep=2pt, text width=1.3cm, minimum height=1.1cm, font=\scriptsize},
    nicht/.style={draw=AlertTitle, very thick, ellipse, align=center,
                  inner sep=2pt, text width=1.3cm, minimum height=1.1cm, font=\scriptsize},
    kante/.style={->, >=Stealth, thick},
    lab/.style={font=\footnotesize\itshape, fill=white, inner sep=1pt}
  ]
    \node[test]  (wo)   at (0,0)       {Wochenende};
    \node[test]  (re)   at (-2.6,-1.8) {Regen};
    \node[nicht] (n1)   at (2.6,-1.8)  {nicht Grillen};
    \node[rad]   (g1)   at (-4.2,-3.6) {Grillen};
    \node[nicht] (n2)   at (-1.0,-3.6) {nicht Grillen};

    \draw[kante] (wo) -- (re) node[lab, midway] {ja};
    \draw[kante] (wo) -- (n1) node[lab, midway] {nein};
    \draw[kante] (re) -- (g1) node[lab, midway] {nein};
    \draw[kante] (re) -- (n2) node[lab, midway] {ja};
  \end{tikzpicture}}
\end{enumerate}
}


% ============================================================
% AUFGABE 2 (schwierig, binäre Attribute)
% ============================================================
\aufgabe{
Gegeben sei der folgende Trainingsdatensatz. Zielattribut ist \textbf{Joggen}.
Alle Attribute sind binär.
\begin{center}
{\scriptsize
\renewcommand{\arraystretch}{1.25}
\setlength{\tabcolsep}{6pt}
\rowcolors{2}{SecondaryColor!12}{white}
\begin{tabular}{c c c c c}
\rowcolor{PrimaryColor}
\textcolor{white}{\textbf{Nr.}} &
\textcolor{white}{\textbf{Regen}} &
\textcolor{white}{\textbf{Windig}} &
\textcolor{white}{\textbf{Kalt}} &
\textcolor{white}{\textbf{Joggen}}\\
1 & Ja   & Nein & Nein & Nein\\
2 & Ja   & Ja   & Ja   & Nein\\
3 & Ja   & Nein & Ja   & Nein\\
4 & Nein & Nein & Nein & Ja\\
5 & Nein & Nein & Ja   & Ja\\
6 & Nein & Nein & Ja   & Ja\\
7 & Nein & Ja   & Nein & Ja\\
8 & Nein & Ja   & Ja   & Nein\\
9 & Nein & Ja   & Ja   & Nein\\
\end{tabular}}
\end{center}

Erstellen Sie mit dem ID3-Algorithmus einen Entscheidungsbaum:
\begin{enumerate}[label=\alph*)]
  \item Bestimmen Sie über die bedingte Entropie $H(V \mid W)$ das Wurzelattribut.
  \item Führen Sie die Attributwahl rekursiv fort und zeichnen Sie den
        vollständigen Entscheidungsbaum.
\end{enumerate}
}

\loesung{
Von $9$ Beispielen sind $4\times$ Joggen ja und $5\times$ nein. Es genügt, die
bedingte Entropie $H(V \mid W)$ zu berechnen und das \textbf{kleinste} zu wählen.

\begin{enumerate}[label=\alph*)]
  \item \textbf{Wurzel:} bedingte Entropie je Attribut.
  \[
    H(V \mid \text{Regen}) = -\Bigl[\,
      \underbrace{\tfrac{3}{9}}_{\text{ja}}\bigl(\tfrac{3}{3}\log_2\tfrac{3}{3}\bigr)
      +
      \underbrace{\tfrac{6}{9}}_{\text{nein}}
      \bigl(\tfrac{4}{6}\log_2\tfrac{4}{6} + \tfrac{2}{6}\log_2\tfrac{2}{6}\bigr)
    \,\Bigr] \approx 0{,}6122
  \]
  \[
    H(V \mid \text{Windig}) = -\Bigl[\,
      \underbrace{\tfrac{5}{9}}_{\text{nein}}
      \bigl(\tfrac{3}{5}\log_2\tfrac{3}{5} + \tfrac{2}{5}\log_2\tfrac{2}{5}\bigr)
      +
      \underbrace{\tfrac{4}{9}}_{\text{ja}}
      \bigl(\tfrac{1}{4}\log_2\tfrac{1}{4} + \tfrac{3}{4}\log_2\tfrac{3}{4}\bigr)
    \,\Bigr] \approx 0{,}9000
  \]
  \[
    H(V \mid \text{Kalt}) = -\Bigl[\,
      \underbrace{\tfrac{3}{9}}_{\text{nein}}
      \bigl(\tfrac{2}{3}\log_2\tfrac{2}{3} + \tfrac{1}{3}\log_2\tfrac{1}{3}\bigr)
      +
      \underbrace{\tfrac{6}{9}}_{\text{ja}}
      \bigl(\tfrac{2}{6}\log_2\tfrac{2}{6} + \tfrac{4}{6}\log_2\tfrac{4}{6}\bigr)
    \,\Bigr] \approx 0{,}9183
  \]
  Kleinste bedingte Entropie: \textbf{Regen} $\Rightarrow$ Wurzelattribut.

  \item \textbf{Rekursion.} Der Zweig \textbf{Regen = ja} (Beispiele 1--3) ist
  rein (nur Joggen nein), also ein Blatt.

  \medskip
  Im Zweig \textbf{Regen = nein} (Beispiele 4--9; $4\times$ ja, $2\times$ nein)
  verbleiben \textbf{Windig} und \textbf{Kalt}:
  \[
    H(V \mid \text{Windig}) = -\Bigl[\,
      \underbrace{\tfrac{3}{6}}_{\text{nein}}\bigl(\tfrac{3}{3}\log_2\tfrac{3}{3}\bigr)
      +
      \underbrace{\tfrac{3}{6}}_{\text{ja}}
      \bigl(\tfrac{1}{3}\log_2\tfrac{1}{3} + \tfrac{2}{3}\log_2\tfrac{2}{3}\bigr)
    \,\Bigr] \approx 0{,}4591
  \]
  \[
    H(V \mid \text{Kalt}) = -\Bigl[\,
      \underbrace{\tfrac{2}{6}}_{\text{nein}}\bigl(\tfrac{2}{2}\log_2\tfrac{2}{2}\bigr)
      +
      \underbrace{\tfrac{4}{6}}_{\text{ja}}
      \bigl(\tfrac{2}{4}\log_2\tfrac{2}{4} + \tfrac{2}{4}\log_2\tfrac{2}{4}\bigr)
    \,\Bigr] \approx 0{,}6667
  \]
  Kleinste bedingte Entropie: \textbf{Windig}. Der Zweig \textbf{Windig = nein}
  ist rein (Joggen). Im Zweig \textbf{Windig = ja} (Beispiele 7--9) trennt das
  letzte Attribut \textbf{Kalt} exakt: nein $\Rightarrow$ Joggen,
  ja $\Rightarrow$ nicht Joggen.

  \vspace{0.4em}
  \centering
  \resizebox{0.85\linewidth}{!}{%
  \begin{tikzpicture}[
    test/.style={draw=PrimaryColor, very thick, rounded corners=4pt,
                 align=center, inner sep=5pt, text width=1.7cm, font=\footnotesize},
    rad/.style={draw=ExampleTitle, very thick, ellipse, align=center,
                inner sep=2pt, text width=1.3cm, minimum height=1.1cm, font=\scriptsize},
    nicht/.style={draw=AlertTitle, very thick, ellipse, align=center,
                  inner sep=2pt, text width=1.3cm, minimum height=1.1cm, font=\scriptsize},
    kante/.style={->, >=Stealth, thick},
    lab/.style={font=\footnotesize\itshape, fill=white, inner sep=1pt}
  ]
    \node[test]  (regen) at (0,0)       {Regen};
    \node[test]  (wind)  at (-3.2,-1.9) {Windig};
    \node[nicht] (n1)    at (3.2,-1.9)  {nicht Joggen};
    \node[rad]   (j1)    at (-5.4,-3.8) {Joggen};
    \node[test]  (kalt)  at (-1.6,-3.8) {Kalt};
    \node[rad]   (j2)    at (-3.2,-5.7) {Joggen};
    \node[nicht] (n2)    at (0.0,-5.7)  {nicht Joggen};

    \draw[kante] (regen) -- (wind) node[lab, midway] {nein};
    \draw[kante] (regen) -- (n1)   node[lab, midway] {ja};
    \draw[kante] (wind)  -- (j1)   node[lab, midway] {nein};
    \draw[kante] (wind)  -- (kalt) node[lab, midway] {ja};
    \draw[kante] (kalt)  -- (j2)   node[lab, midway] {nein};
    \draw[kante] (kalt)  -- (n2)   node[lab, midway] {ja};
  \end{tikzpicture}}
\end{enumerate}
}